\section{Axiomatic Foundations}
\label{sec:axioms}

The five theoretical pillars of Section~\ref{sec:foundations} provide the conceptual vocabulary for the \Apeiron{} Framework. To operationalize these pillars into a formal architecture, we posit twelve axioms. These axioms are not arbitrary stipulations; they encode the ontological and epistemological commitments that distinguish the framework from conventional AI. Each axiom is referenced in the layer specifications that follow, grounding the architectural choices in a rigorous foundation.

\begin{axiom}[Irreducibility]\label{ax:irred}
There exists a set $\Obs$ of \emph{foundational observables}. For each $o \in \Obs$, there is no non-trivial decomposition $\decomp(o)$ that preserves identity and generative capacity.
\end{axiom}

\begin{axiom}[Multi-Axial Atomicity]\label{ax:multi}
The atomicity of an observable is validated by convergence across at least three independent mathematical frameworks (Boolean, measure, categorical, information).
\end{axiom}

\begin{axiom}[Relational Constitution]\label{ax:relational}
The identity of any entity at layer $L_k$ ($k \ge 2$) is constituted by its relational position within a hypergraph of lower-layer entities.
\end{axiom}

\begin{axiom}[Emergence]\label{ax:emergence}
For each layer $L_{k+1}$, there exist properties and behaviors that are not reducible to the sum of properties of entities at layer $L_k$.
\end{axiom}

\begin{axiom}[Temporal Endogeneity]\label{ax:time}
Time is not a primitive background but an emergent ordering derived from the causal structure of state transitions.
\end{axiom}

\begin{axiom}[Ontological Openness]\label{ax:openness}
The set of categories and concepts available to the system is not fixed but can be extended through ontogenesis (Layer~13).
\end{axiom}

\begin{axiom}[Autopoietic Closure]\label{ax:autopoiesis}
A world-model $W$ is autopoietic if it contains a sub-network of processes that regenerate all components of $W$.
\end{axiom}

\begin{axiom}[Ethical Reversibility]\label{ax:reversibility}
Any morphogenetic action at Layer~15 must be designed with a verifiable rollback mechanism, unless a supermajority of a designated governance council approves irreversibility.
\end{axiom}

\begin{axiom}[Observer Relativity]\label{ax:observer}
All atomicity judgments are relative to a \texttt{MetaSpecification} that encodes the observer's framework.
\end{axiom}

\begin{axiom}[Generativity]\label{ax:generativity}
An observable's capacity to serve as a substrate for higher layers is proportional to its resonance richness, relational degree, and atomicity consensus.
\end{axiom}

\begin{axiom}[Epistemic Humility]\label{ax:humility}
The system must maintain explicit representations of the uncertainty and provenance of all meta-decisions.
\end{axiom}

\begin{axiom}[Convergence]\label{ax:convergence}
The framework asymptotically approaches a fixed point of self-consistency (Layer~17 equilibrium) under iterative refinement.
\end{axiom}

These axioms are not isolated declarations; they form a coherent system. Axioms~\ref{ax:irred}--\ref{ax:multi} ground the substrate. Axioms~\ref{ax:relational}--\ref{ax:emergence} govern emergence. Axioms~\ref{ax:time}--\ref{ax:openness} enable ontological fluidity. Axioms~\ref{ax:autopoiesis}--\ref{ax:reversibility} embed ethics and world-making. Axioms~\ref{ax:observer}--\ref{ax:convergence} ensure reflexivity and asymptotic self-consistency. In the layer descriptions that follow, each layer explicitly states which axioms it satisfies.